% در این فایل، عنوان پایان‌نامه، مشخصات خود و چکیده پایان‌نامه را به انگلیسی، وارد کنید.
% توجه داشته باشید که جدول حاوی مشخصات پایان‌نامه/رساله، به طور خودکار، رسم می‌شود.
%%%%%%%%%%%%%%%%%%%%%%%%%%%%%%%%%%%%
\baselineskip=.6cm
\begin{latin}
    \latinfaculty{Faculty of Mathematical Sciences}
    \latindegree{Master of Computer Science}
    \latinuniversity{Tarbiat Modares University}
% group
\latinsubject{Department of Computer Science}
\latinfield{Specialization in Data Mining}
\latintitle{Improving Diffusion Models Based on Vision Transformers for Image Generation}
\firstlatinsupervisor{Dr.\ Mansoor Rezghi Ahaq}
%\secondlatinsupervisor{Second Supervisor}
%\firstlatinadvisor{First Advisor}
%\secondlatinadvisor{Second Advisor}
\latinname{Mostafa}
\latinsurname{Shahbazi Dil}
\latinthesisdate{2026 Feb}
\latinkeywords{Vision transformers, diffusion models, local and global features, feature aggregation, parameter-efficient fine-tuning}
\newpage
\thispagestyle{empty}
\baselineskip=.750cm
\en-abstract{%\noindent
Diffusion transformer models, such as DiT, provide high quality in class-conditional image generation, but fully fine-tuning them for new domains or requirements is costly in terms of memory, training time, and storage. This thesis introduces DuoDiT as a parameter-efficient fine-tuning architecture for DiT models. In this architecture, the main DiT-XL/2 backbone remains frozen, while a second branch with a smaller patch size extracts local and fine-grained features. To align this branch with the backbone tokens, DuoDiT uses a local class-token mechanism that summarizes each spatial group of fine-grained tokens in a content-aware manner. The main evaluation is performed on 50,000 images generated conditionally on the ImageNet-1K dataset at $256 \times 256$ resolution. DuoDiT-200, with 17.95M trainable parameters out of 690.39M total parameters, achieves FID=2.219, KID=0.0014, and LPIPS-Mean=0.7162. In addition, a human study shows that DuoDiT samples are preferred over LightningDiT in 53.26\% of the comparisons. The results indicate that adding a lightweight high-resolution pathway can improve perceptual quality and fine-tuning efficiency without fully retraining the model.
}

%\markboth{\lr{Abstract}}{\lr{Abstract}}
\noindent
\begin{large}
\begin{center}
Abstract
\end{center}
\end{large}
%
Diffusion transformer models, such as DiT, provide high quality in class-conditional image generation, but fully fine-tuning them for new domains or requirements is costly in terms of memory, training time, and storage. This thesis introduces DuoDiT as a parameter-efficient fine-tuning architecture for DiT models. In this architecture, the main DiT-XL/2 backbone remains frozen, while a second branch with a smaller patch size extracts local and fine-grained features. To align this branch with the backbone tokens, DuoDiT uses a local class-token mechanism that summarizes each spatial group of fine-grained tokens in a content-aware manner. The main evaluation is performed on 50,000 images generated conditionally on the ImageNet-1K dataset at $256 \times 256$ resolution. DuoDiT-200, with 17.95M trainable parameters out of 690.39M total parameters, achieves FID=2.219, KID=0.0014, and LPIPS-Mean=0.7162. In addition, a human study shows that DuoDiT samples are preferred over LightningDiT in 53.26\% of the comparisons. The results indicate that adding a lightweight high-resolution pathway can improve perceptual quality and fine-tuning efficiency without fully retraining the model.
\newline
{{\bf Keywords:}}
Vision transformers, diffusion models, local and global features, feature aggregation, parameter-efficient fine-tuning%}
\newpage
\latinvtitle%
%\thispagestyle{empty}
\end{latin}
